\section{Quasi-compact spaces and maps}
\label{section-quasi-compact}

\noindent
The phrase ``compact'' will be reserved
for Hausdorff topological spaces. And many spaces occurring
in algebraic geometry are not Hausdorff.

\begin{definition}
\label{definition-quasi-compact}
Quasi-compactness.
\begin{enumerate}
\item We say that a topological space $X$ is {\it quasi-compact}
if every open covering of $X$ has a finite subcover.
\item We say that a continuous map $f : X \to Y$ is {\it quasi-compact}
if the inverse image $f^{-1}(V)$ of every quasi-compact open $V \subset Y$
is quasi-compact.
\item We say a subset $Z \subset X$ is {\it retrocompact}
if the inclusion map $Z \to X$ is quasi-compact.
\end{enumerate}
\end{definition}

\noindent
In many texts on topology a space is called {\it compact} if it
is quasi-compact and Hausdorff; and in other texts the Hausdorff
condition is omitted. To avoid confusion in algebraic geometry
we use the term quasi-compact. The notion of quasi-compactness
of a map is very different from the notion of a ``proper map'',
since there we require (besides closedness and separatedness)
the inverse image of any quasi-compact subset of the target to be
quasi-compact, whereas in the definition above we only consider
quasi-compact {\it open} sets.

\begin{lemma}
\label{lemma-composition-quasi-compact}
A composition of quasi-compact maps is quasi-compact.
\end{lemma}

\begin{proof}
This is immediate from the definition.
\end{proof}

\begin{lemma}
\label{lemma-closed-in-quasi-compact}
A closed subset of a quasi-compact topological space is quasi-compact.
\end{lemma}

\begin{proof}
Let $E \subset X$ be a closed subset of the quasi-compact space $X$.
Let $E = \bigcup V_j$ be an open covering. Choose $U_j \subset X$
open such that $V_j = E \cap U_j$. Then $X = (X \setminus E) \cup \bigcup U_j$
is an open covering of $X$. Hence
$X = (X \setminus E) \cup U_{j_1} \cup \ldots \cup U_{j_n}$ for some
$n$ and indices $j_i$. Thus $E = V_{j_1} \cup \ldots \cup V_{j_n}$
as desired.
\end{proof}

\begin{lemma}
\label{lemma-quasi-compact-in-Hausdorff}
Let $X$ be a Hausdorff topological space.
\begin{enumerate}
\item If $E \subset X$ is quasi-compact, then it is closed.
\item If $E_1, E_2 \subset X$ are disjoint quasi-compact subsets
then there exists opens $E_i \subset U_i$ with $U_1 \cap U_2 = \emptyset$.
\end{enumerate}
\end{lemma}

\begin{proof}
Proof of (1). Let $x \in X$, $x \not \in E$.
For every $e \in E$ we can find disjoint opens $V_e$ and $U_e$
with $e \in V_e$ and $x \in U_e$. Since $E \subset \bigcup V_e$
we can find finitely many $e_1, \ldots, e_n$ such that
$E \subset V_{e_1} \cup \ldots \cup V_{e_n}$. Then
$U = U_{e_1} \cap \ldots \cap U_{e_n}$ is an open neighbourhood
of $x$ which avoids $V_{e_1} \cup \ldots \cup V_{e_n}$. In particular
it avoids $E$. Thus $E$ is closed.

\medskip\noindent
Proof of (2). In the proof of (1) we have seen that given $x \in E_1$
we can find an open neighbourhood $x \in U_x$ and an open
$E_2 \subset V_x$ such that $U_x \cap V_x = \emptyset$. Because
$E_1$ is quasi-compact we can find a finite number $x_i \in E_1$
such that $E_1 \subset U = U_{x_1} \cup \ldots \cup U_{x_n}$.
We take $V = V_{x_1} \cap \ldots \cap V_{x_n}$ to finish the proof.
\end{proof}

\begin{lemma}
\label{lemma-closed-in-compact}
Let $X$ be a quasi-compact Hausdorff space. Let $E \subset X$.
The following are equivalent: (a) $E$ is closed in $X$, (b)
$E$ is quasi-compact.
\end{lemma}

\begin{proof}
The implication (a) $\Rightarrow$ (b) is
Lemma \ref{lemma-closed-in-quasi-compact}.
The implication (b) $\Rightarrow$ (a) is
Lemma \ref{lemma-quasi-compact-in-Hausdorff}.
\end{proof}

\noindent
The following is really a reformulation of the
quasi-compact property.

\begin{lemma}
\label{lemma-intersection-closed-in-quasi-compact}
Let $X$ be a quasi-compact topological space.
If $\{Z_\alpha\}_{\alpha \in A}$ is a collection of closed subsets
such that the intersection of each finite subcollection
is nonempty, then $\bigcap_{\alpha \in A} Z_\alpha$ is nonempty.
\end{lemma}

\begin{proof}
We suppose that $\bigcap_{\alpha \in A} Z_{\alpha} = \emptyset$.
So we have $\bigcup_{\alpha \in A} (X\setminus Z_{\alpha})=X$ by
complementation. As the subsets $Z_\alpha$ are closed,
$\bigcup_{\alpha \in A} (X \setminus Z_{\alpha})$
is an open covering of the quasi-compact space $X$.
Thus there exists a finite subset $J \subset A$ such that
$X = \bigcup_{\alpha \in J} (X\setminus Z_{\alpha})$.
The complementary is then empty, which means that
$\bigcap_{\alpha \in J} Z_{\alpha} = \emptyset$. It proves there exists a
finite subcollection of $\{Z_{\alpha}\}_{\alpha \in J}$ verifying
$\bigcap_{\alpha \in J} Z_{\alpha}=\emptyset$, which concludes by
contraposition.
\end{proof}

\begin{lemma}
\label{lemma-image-quasi-compact}
Let $f : X \to Y$ be a continuous map of topological spaces.
\begin{enumerate}
\item If $X$ is quasi-compact, then $f(X)$ is quasi-compact.
\item If $f$ is quasi-compact, then $f(X)$ is retrocompact.
\end{enumerate}
\end{lemma}

\begin{proof}
If $f(X) = \bigcup V_i$ is an open covering, then $X = \bigcup f^{-1}(V_i)$
is an open covering. Hence if $X$ is quasi-compact then
$X = f^{-1}(V_{i_1}) \cup \ldots \cup f^{-1}(V_{i_n})$ for some
$i_1, \ldots, i_n \in I$ and hence
$f(X) = V_{i_1} \cup \ldots \cup V_{i_n}$. This proves (1).
Assume $f$ is quasi-compact, and let $V \subset Y$ be quasi-compact open.
Then $f^{-1}(V)$ is quasi-compact, hence by (1) we see that
$f(f^{-1}(V)) = f(X) \cap V$ is quasi-compact. Hence $f(X)$
is retrocompact.
\end{proof}

\begin{lemma}
\label{lemma-quasi-compact-closed-point}
Let $X$ be a topological space. Assume that
\begin{enumerate}
\item $X$ is nonempty,
\item $X$ is quasi-compact, and
\item $X$ is Kolmogorov.
\end{enumerate}
Then $X$ has a closed point.
\end{lemma}

\begin{proof}
Consider the set
$$
\mathcal{T} =
\{Z \subset X \mid Z = \overline{\{x\}} \text{ for some }x \in X\}
$$
of all closures of singletons in $X$. It is nonempty since $X$ is
nonempty. Make $\mathcal{T}$ into a
partially ordered set using the relation of inclusion.
Suppose $Z_\alpha$, $\alpha \in A$ is a totally ordered subset of $\mathcal{T}$.
By Lemma \ref{lemma-intersection-closed-in-quasi-compact} we see
that $\bigcap_{\alpha \in A} Z_\alpha \not = \emptyset$. Hence there exists
some $x \in \bigcap_{\alpha \in A} Z_\alpha$ and we see that
$Z = \overline{\{x\}}\in \mathcal{T}$ is a lower bound for
the family. By Zorn's lemma there exists a minimal element
$Z \in \mathcal{T}$. As $X$ is Kolmogorov we conclude that
$Z = \{x\}$ for some $x$ and $x \in X$ is a closed point.
\end{proof}

\begin{lemma}
\label{lemma-closed-points-quasi-compact}
Let $X$ be a quasi-compact Kolmogorov space. Then the set $X_0$ of
closed points of $X$ is quasi-compact.
\end{lemma}

\begin{proof}
Let $X_0 = \bigcup U_{i, 0}$ be an open covering.
Write $U_{i, 0} = X_0 \cap U_i$ for some open $U_i \subset X$.
Consider the complement $Z$ of $\bigcup U_i$. This is a closed subset of
$X$, hence quasi-compact (Lemma \ref{lemma-closed-in-quasi-compact})
and Kolmogorov. By Lemma \ref{lemma-quasi-compact-closed-point}
if $Z$ is nonempty it would have a closed
point which contradicts the fact that $X_0 \subset \bigcup U_i$.
Hence $Z = \emptyset$ and $X = \bigcup U_i$. Since $X$ is quasi-compact
this covering has a finite subcover and we conclude.
\end{proof}

\begin{lemma}
\label{lemma-connected-component-intersection}
Let $X$ be a topological space.
Assume
\begin{enumerate}
\item $X$ is quasi-compact,
\item $X$ has a basis for the topology consisting of quasi-compact opens, and
\item the intersection of two quasi-compact opens is quasi-compact.
\end{enumerate}
For any $x \in X$ the connected component of $X$ containing
$x$ is the intersection of all open and closed subsets
of $X$ containing $x$.
\end{lemma}

\begin{proof}
Let $T$ be the connected component containing $x$.
Let $S = \bigcap_{\alpha \in A} Z_\alpha$ be the intersection of all
open and closed subsets $Z_\alpha$ of $X$ containing $x$.
Note that $S$ is closed in $X$.
Note that any finite intersection of $Z_\alpha$'s is a $Z_\alpha$.
Because $T$ is connected and $x \in T$ we have $T \subset S$.
It suffices to show that $S$ is connected.
If not, then there exists a disjoint union decomposition
$S = B \amalg C$ with $B$ and $C$ open and closed in $S$.
In particular, $B$ and $C$ are closed in $X$, and so quasi-compact by
Lemma \ref{lemma-closed-in-quasi-compact} and assumption (1).
By assumption (2) there exist quasi-compact opens
$U, V \subset X$ with $B = S \cap U$ and $C = S \cap V$ (details omitted).
Then $U \cap V \cap S = \emptyset$.
Hence $\bigcap_\alpha U \cap V \cap Z_\alpha = \emptyset$.
By assumption (3) the intersection $U \cap V$ is quasi-compact.
By Lemma \ref{lemma-intersection-closed-in-quasi-compact}
for some $\alpha' \in A$ we have $U \cap V \cap Z_{\alpha'} = \emptyset$.
Since $X \setminus (U \cup V)$ is disjoint from $S$
and closed in $X$ hence quasi-compact, we can use the same lemma
to see that $Z_{\alpha''} \subset U \cup V$ for some $\alpha'' \in A$.
Then $Z_\alpha = Z_{\alpha'} \cap Z_{\alpha''}$ is contained
in $U \cup V$ and disjoint from $U \cap V$.
Hence $Z_\alpha = U \cap Z_\alpha \amalg V \cap Z_\alpha$
is a decomposition into two open pieces,
hence $U \cap Z_\alpha$ and $V \cap Z_\alpha$ are open and closed in $X$.
Thus, if $x \in B$ say, then we see that $S \subset U \cap Z_\alpha$
and we conclude that $C = \emptyset$.
\end{proof}

\begin{lemma}
\label{lemma-connected-component-intersection-compact-Hausdorff}
Let $X$ be a topological space. Assume $X$ is quasi-compact and Hausdorff.
For any $x \in X$ the connected component of $X$ containing
$x$ is the intersection of all open and closed subsets
of $X$ containing $x$.
\end{lemma}

\begin{proof}
Let $T$ be the connected component containing $x$.
Let $S = \bigcap_{\alpha \in A} Z_\alpha$ be the intersection of all
open and closed subsets $Z_\alpha$ of $X$ containing $x$.
Note that $S$ is closed in $X$.
Note that any finite intersection of $Z_\alpha$'s is a $Z_\alpha$.
Because $T$ is connected and $x \in T$ we have $T \subset S$.
It suffices to show that $S$ is connected.
If not, then there exists a disjoint union decomposition
$S = B \amalg C$ with $B$ and $C$ open and closed in $S$.
In particular, $B$ and $C$ are closed in $X$, and so quasi-compact by
Lemma \ref{lemma-closed-in-quasi-compact}.
By Lemma \ref{lemma-quasi-compact-in-Hausdorff}
there exist disjoint opens $U, V \subset X$ with $B \subset U$ and
$C \subset V$. Then $X \setminus U \cup V$ is closed in $X$
hence quasi-compact (Lemma \ref{lemma-closed-in-quasi-compact}).
It follows that $(X \setminus U \cup V) \cap Z_\alpha = \emptyset$
for some $\alpha$ by Lemma \ref{lemma-intersection-closed-in-quasi-compact}.
In other words, $Z_\alpha \subset U \cup V$. Thus
$Z_\alpha = Z_\alpha \cap V \amalg Z_\alpha \cap U$
is a decomposition into two open pieces,
hence $U \cap Z_\alpha$ and $V \cap Z_\alpha$ are open and closed in $X$.
Thus, if $x \in B$ say, then we see that $S \subset U \cap Z_\alpha$
and we conclude that $C = \emptyset$.
\end{proof}

\begin{lemma}
\label{lemma-closed-union-connected-components}
Let $X$ be a topological space.
Assume
\begin{enumerate}
\item $X$ is quasi-compact,
\item $X$ has a basis for the topology consisting of quasi-compact opens, and
\item the intersection of two quasi-compact opens is quasi-compact.
\end{enumerate}
For a subset $T \subset X$ the following are equivalent:
\begin{enumerate}
\item[(a)] $T$ is an intersection of open and closed subsets of $X$, and
\item[(b)] $T$ is closed in $X$ and is a union of connected components of $X$.
\end{enumerate}
\end{lemma}

\begin{proof}
It is clear that (a) implies (b).
Assume (b). Let $x \in X$, $x \not \in T$. Let $x \in C \subset X$
be the connected component of $X$ containing $x$. By
Lemma \ref{lemma-connected-component-intersection}
we see that $C = \bigcap V_\alpha$ is the intersection of all open and
closed subsets $V_\alpha$ of $X$ which contain $C$.
In particular, any pairwise intersection $V_\alpha \cap V_\beta$
occurs as a $V_\alpha$.
As $T$ is a union of connected components
of $X$ we see that $C \cap T = \emptyset$. Hence
$T \cap \bigcap V_\alpha = \emptyset$. Since $T$ is quasi-compact as a
closed subset of a quasi-compact space (see
Lemma \ref{lemma-closed-in-quasi-compact})
we deduce that $T \cap V_\alpha = \emptyset$ for some $\alpha$, see
Lemma \ref{lemma-intersection-closed-in-quasi-compact}.
For this $\alpha$ we see that $U_\alpha = X \setminus V_\alpha$
is an open and closed subset of $X$ which contains $T$ and not $x$.
The lemma follows.
\end{proof}

\begin{lemma}
\label{lemma-Noetherian-quasi-compact}
Let $X$ be a Noetherian topological space.
\begin{enumerate}
\item The space $X$ is quasi-compact.
\item Any subset of $X$ is retrocompact.
\end{enumerate}
\end{lemma}

\begin{proof}
Suppose $X = \bigcup U_i$ is an open covering of $X$ indexed
by the set $I$ which does not have a refinement by a finite
open covering. Choose $i_1, i_2, \ldots $ elements of $I$ inductively
in the following way: Choose $i_{n + 1}$ such that $U_{i_{n + 1}}$
is not contained in $U_{i_1} \cup \ldots \cup U_{i_n}$. Thus we see that
$X \supset (X \setminus U_{i_1}) \supset
(X \setminus U_{i_1} \cup U_{i_2}) \supset \ldots$ is a strictly
decreasing infinite sequence of closed subsets. This contradicts
the fact that $X$ is Noetherian. This proves the first assertion.
The second assertion is now clear since every subset of $X$ is Noetherian by
Lemma \ref{lemma-Noetherian}.
\end{proof}

\begin{lemma}
\label{lemma-quasi-compact-locally-Noetherian-Noetherian}
A quasi-compact locally Noetherian space is Noetherian.
\end{lemma}

\begin{proof}
The conditions imply immediately that $X$ has a finite covering by
Noetherian subsets, and hence is Noetherian by
Lemma \ref{lemma-finite-union-Noetherian}.
\end{proof}

\begin{lemma}[Alexander subbase theorem]
\label{lemma-subbase-theorem}
Let $X$ be a topological space. Let $\mathcal{B}$ be a subbase for $X$.
If every covering of $X$ by elements of $\mathcal{B}$ has a finite
refinement, then $X$ is quasi-compact.
\end{lemma}

\begin{proof}
Assume there is an open covering of $X$ which does not have a finite
refinement. Using Zorn's lemma we can choose a maximal open covering
$X = \bigcup_{i \in I} U_i$ which does not have a finite refinement
(details omitted).
In other words, if $U \subset X$ is any open which does not occur as
one of the $U_i$, then the covering $X = U \cup \bigcup_{i \in I} U_i$
does have a finite refinement. Let $I' \subset I$ be the set of indices
such that $U_i \in \mathcal{B}$. Then $\bigcup_{i \in I'} U_i \not = X$,
since otherwise we would get a finite refinement covering $X$ by our
assumption on $\mathcal{B}$. Pick $x \in X$,
$x \not \in \bigcup_{i \in I'} U_i$. Pick $i \in I$ with $x \in U_i$.
Pick $V_1, \ldots, V_n \in \mathcal{B}$ such that
$x \in V_1 \cap \ldots \cap V_n \subset U_i$. This is
possible as $\mathcal{B}$ is a subbasis for $X$. Note that
$V_j$ does not occur as a $U_i$. By maximality of the chosen
covering we see that for each $j$ there exist
$i_{j, 1}, \ldots, i_{j, n_j} \in I$ such that
$X = V_j \cup U_{i_{j, 1}} \cup \ldots \cup U_{i_{j, n_j}}$.
Since $V_1 \cap \ldots \cap V_n \subset U_i$ we conclude that
$X = U_i \cup \bigcup U_{i_{j, l}}$ a contradiction.
\end{proof}
